% =====================================================================
%  DBMS Lab Record - Section D / 6. Project (Workout Tracker)
%  Continuation of the record: page numbering starts after 120.
%  Build: latexmk -xelatex section-d-project.tex
% =====================================================================
% =====================================================================
%  DBMS Lab Record Book
%  Build: latexmk -pdf record.tex
% =====================================================================
\documentclass[12pt,a4paper]{article}

\usepackage{fontspec}
\setmainfont{Liberation Serif}         % metric-compatible with Times New Roman
\setmonofont{Nimbus Mono PS}[Scale=0.87]
\usepackage[a4paper,left=0.75in,right=0.75in,top=1in,bottom=1in]{geometry}
\usepackage{tikz}
\usepackage{eso-pic}
\usepackage{tabularx}
\usepackage{array}
\usepackage{fancyvrb}
\usepackage[normalem]{ulem}

\pagestyle{empty}
\setlength{\parindent}{0pt}
\setlength{\parskip}{0pt}

% ---------- page border (single line, 30pt from the page edge) ----------
\AddToShipoutPictureBG{%
  \begin{tikzpicture}[remember picture,overlay]
    \draw[line width=0.5pt]
      ([xshift=30pt,yshift=-30pt]current page.north west)
      rectangle
      ([xshift=-30pt,yshift=30pt]current page.south east);
  \end{tikzpicture}%
}

% ---------- page number: outside the border, 13pt from the page edge ------
\AddToShipoutPictureFG{%
  \begin{tikzpicture}[remember picture,overlay]
    \node[anchor=south, inner sep=0pt] at ([yshift=13pt]current page.south)
      {\fontsize{12}{14}\selectfont\thepage};
  \end{tikzpicture}%
}

% ---------- headings (sizes taken from DBMS_lab_Format.docx) ----------
\newcommand{\mainheading}[1]{%
  \begin{center}\textbf{\fontsize{18}{22}\selectfont #1}\end{center}\vspace{6pt}}
\newcommand{\subheading}[1]{%
  \textbf{\fontsize{16}{19}\selectfont\uline{#1}}\par\vspace{10pt}}
\newcommand{\subsubheading}[1]{%
  \textbf{\fontsize{14}{17}\selectfont\uline{#1}}\par\vspace{10pt}}
\newcommand{\qheading}[1]{%
  \textbf{\fontsize{14}{17}\selectfont\uline{#1}}\par\vspace{10pt}}
\newcommand{\codeheading}[1]{%
  \textbf{\fontsize{12}{15}\selectfont\uline{#1}}\par\vspace{4pt}}

% ---------- code / output blocks ----------
\DefineVerbatimEnvironment{sqlcode}{Verbatim}%
  {fontsize=\small,formatcom=\ttfamily,xleftmargin=0pt}
\DefineVerbatimEnvironment{sqlout}{Verbatim}%
  {fontsize=\small,formatcom=\ttfamily,xleftmargin=0pt}


\usepackage{graphicx}
\graphicspath{{img/}}

\setcounter{page}{121}

% ---------- report headings (layout taken from the sample report) ----------
\newcommand{\projtitle}[1]{%
  \textbf{\fontsize{20}{24}\selectfont #1}\par\vspace{14pt}}
\newcommand{\projhead}[1]{%
  \textbf{\fontsize{15}{18}\selectfont #1}\par\vspace{8pt}}
\newcommand{\projsub}[1]{%
  \textbf{\fontsize{13}{16}\selectfont #1}\par\vspace{6pt}}
\newcommand{\projitem}[2]{%
  \noindent\textbf{\fontsize{12.5}{15}\selectfont #1}\ \ {\fontsize{12.5}{15}\selectfont #2}\par\vspace{6pt}}
\newcommand{\projtext}[1]{%
  {\fontsize{12}{15.6}\selectfont #1\par}\vspace{4pt}}
\newcommand{\shot}[2][430pt]{%
  \noindent\includegraphics[width=#1]{#2}\par\vspace{16pt}}

\begin{document}

\mainheading{Section D}
\subheading{6.\ \ Project}

\projtitle{Workout Tracker}

\projhead{Abstract}

{\fontsize{12}{15.6}\selectfont
\setlength{\parskip}{8pt}

The Workout Tracker is a simple Microsoft Access based database application
designed to record a person's gym workouts and the changes in their body over
time. Every training session is stored with its date, start time, duration, the
workout performed, the details of the session and the distance covered, so that
the user can look back at what was done on any day instead of relying on a
notebook.

The database consists of four related tables --- Users, BodyHistory,
WorkoutLog and WorkoutNameMaint. The Users table holds the personal details of
each person, the WorkoutNameMaint table holds the list of workouts such as Abs
Day 1, Chest/Tris TT or DDP Energy, and the WorkoutLog and BodyHistory tables
record one row for each workout session and each body measurement for a user.
The tables are related so that a workout or a measurement always belongs to a
valid user and a log entry always refers to a workout that exists in the list.

Records are entered through a set of user friendly forms rather than by typing
directly into the tables. The main menu opens the workout log form, the body
history form, the user and workout maintenance forms and the reports form, and
each form saves its own record. A set of queries --- including body history
selection, the export queries and the update queries that calculate the energy
burnt by a workout --- supplies the figures that the reports form needs. The
reports form takes a start date, an end date and a workout name and exports the
selected workout records for further analysis.

Overall, the project demonstrates the practical use of Microsoft Access for
building a small scale information system that is easy to use, keeps the data
consistent, and answers common fitness questions without any manual
calculation.

\par}

\clearpage
\projhead{Screenshots}

\projsub{Tables}

\projitem{1.}{Body History Table}
{\small The BodyHistory table stores one row for every measurement taken, with
the user, the entry date, the height, the weight and the body fat percentage.}
\vspace{4pt}

\shot[430pt]{w_bodyhistory_table.png}

\projitem{2.}{Users Table}
{\small The Users table stores the personal details of each person using the
application --- name, address, city, zip code, date of birth and phone number.}
\vspace{4pt}

\shot[430pt]{w_users_table.png}

\clearpage
\projitem{3.}{Workout Name Table}
{\small The WorkoutNameMaint table holds the master list of workout names that
the log and the reports pick from.}
\vspace{4pt}

\shot[430pt]{w_workoutname_table.png}

\vspace{6pt}
\projhead{Forms}

\projitem{1.}{Main Menu}
{\small The main menu is the starting point of the application. It opens the
workout log, the body history, the user and workout maintenance screens and the
reports form from a single window.}
\vspace{4pt}

\shot[430pt]{w_main_menu.png}

\clearpage
\projitem{2.}{Workout Log Form}
{\small The workout log form records a training session. The user is picked
from a list, and the date, start time, duration, workout name, workout details
and distance are stored; the Update Durations button recalculates the duration
of the session.}
\vspace{4pt}

\shot[430pt]{w_workout_log.png}

\projitem{3.}{Body History Form}
{\small The body history form records the height, weight and body fat percentage
measured on a given date, so that the progress of the user can be followed over
time.}
\vspace{4pt}

\shot[430pt]{w_body_history.png}

\clearpage
\projitem{4.}{Workout Maintenance Form}
{\small The workout maintenance form is used to add, edit and delete the workout
names used by the log.}
\vspace{4pt}

\shot[415pt]{w_workout_maint.png}

\vspace{2pt}
\projhead{Reports}

\projitem{1.}{Reports Form}
{\small The reports form asks for a start date, an end date and a workout name,
and exports the matching workout records so that they can be analysed outside
the database.}
\vspace{4pt}

\shot[415pt]{w_reports.png}

\end{document}
